\begin{claim}{D.75}{Proved}{The full negative torus cusp centralizer and unique transporters}
Let $b$ be any negative involution of the full
once-punctured-torus loop-outer group, with the actual
cusp-normalized extension of D.58--D.59:
\[
 \begin{gathered}
 F=\widehat F(x,y),\qquad c=[x,y],\qquad
 E=F\rtimes\langle v\rangle,\\
 v^2=1,\qquad vcv^{-1}=c^{-1},\qquad
 \alpha=\operatorname{Inn}(v)|_F .
 \end{gathered}                                             \tag{D.75.1}
\]
Then
\[
                    C_F(v)=1,\qquad N_E(\langle v\rangle)
                                      =\langle v\rangle .
                                                               \tag{D.75.2}
\]
Every cusp reflection $c^t v$, $t\in\widehat{\mathbb Z}$,
has the corresponding full trivial fixed subgroup in $F$.
In particular there is a unique $g\in F$ with
\[
                            gvg^{-1}=cv .                   \tag{D.75.3}
\]
Its homology is the primitive generator of the full fixed
line specified in D.59.

For every full unit $\mu\in\widehat{\mathbb Z}^{\times}$
and both literal reflection labels $v_i=c^i v$, $i=0,1$,
the entire transporter set is a singleton:
\[
 \begin{gathered}
 \{a\in F:[a,v_i]=c^\mu\}=\{a_{\mu,i}\},\\
 a_{\mu,0}=c^{(\mu-1)/2}g,\qquad
 a_{\mu,1}=c^{(\mu+1)/2}g^{-1}.
 \end{gathered}                                             \tag{D.75.4}
\]
If $\alpha_*\bmod2\ne I$, the extension has one section
class and is real projective.  If $\alpha_*\bmod2=I$,
the other section class retains the fixed group and
surjective homology map of D.59; no full procyclic
identification of that group is asserted here.

With the standard model selected by these section counts,
the complete D.61 peripheral-conjugator existence problem
is exactly
\[
 \text{there exist }\mu\in\widehat{\mathbb Z}^{\times},
 \ i\in\{0,1\}\text{ with }\quad
                     c^{-\mu}a_{\mu,i}\sim_F a_{\mu,i}.
                                                               \tag{D.75.5}
\]
The maps proving this equivalence are retained below.
Existence in (D.75.5), and recognition of an arbitrary
negative action by one of the standard models, remain
unproved.
\end{claim}
\begin{proof}
\emph{The one actual cusp module.}
Let $u:E\to C_2=\mathbb F_2$ be the quotient character,
$u(v)=1$, and set
$L=\overline{\langle c\rangle}$ and $D=N_E(L)$.
D.58 gives the actual inverse homomorphisms
\[
 \begin{aligned}
 \widehat{\mathbb Z}\rtimes_{-1}C_2
        &\longrightarrow D,
        &(t,\tau^\epsilon)&\longmapsto c^t v^\epsilon,\\
 D&\longrightarrow\widehat{\mathbb Z}\rtimes_{-1}C_2,
        &d&\longmapsto(t,u(d)),
        \qquad d v^{-u(d)}=c^t .
 \end{aligned}                                             \tag{D.75.6}
\]
Its two internal section classes have representatives
$S=\langle v\rangle$ and $S'=\langle cv\rangle$.
Their normalizers in $D$ are $S$ and $S'$:
an element $c^a$ commuting with a reflection has
$2a=0$, so $a=0$.  D.58 supplies an actual
$g\in F$ with $S'=gSg^{-1}$ and (D.75.3).
We will prove its uniqueness after proving (D.75.2).

Put $k=\mathbb F_2$, and form the actual compact module
\[
 0\longrightarrow\Delta
 \longrightarrow k[[E/D]]
            \xrightarrow{\ \varepsilon\ }k\longrightarrow0 .
                                                               \tag{D.75.7}
\]
The explicit restriction map
\[
            F/L\longrightarrow E/D,\qquad fL\longmapsto fD
                                                               \tag{D.75.8}
\]
is a homeomorphism: $E=FD$ and $F\cap D=L$
give bijectivity, and both spaces are compact Hausdorff.
In these coordinates the full $E$-action is
\[
                   e\cdot(fL)=efv^{-u(e)}L .              \tag{D.75.9}
\]
The representative is in $F$, and its coset is
independent of the representative of $fL$, because
$v$ normalizes $L$.  Thus $\Delta|_F$ is exactly
the augmentation module of the full completed
once-punctured-torus pair $(F,L)$, with the original
commutator cusp and the transported $E$-action.

For a finite continuous $E$-module $M$ over $k$, define
\[
 \mathscr H_M=\operatorname{Hom}_{k,\mathrm{cont}}(\Delta,M),
 \qquad (e\varphi)(d)=e\varphi(e^{-1}d).
                                                               \tag{D.75.10}
\]
This is the same actual discrete coefficient
construction as (D.74.9), now for (D.75.7).
The augmentation sequence, its continuous Hom sequence,
and Shapiro give
\[
 \begin{gathered}
 R^n(E,M)=H^n(E,D;M)=H^{n-1}(E,\mathscr H_M)
                       \quad(n\geq1),\\
 \cdots\longrightarrow H^n(E,M)\xrightarrow{\ r\ }H^n(D,M)
 \longrightarrow R^{n+1}(E,M)
 \longrightarrow H^{n+1}(E,M)\longrightarrow\cdots .
 \end{gathered}                                             \tag{D.75.11}
\]
All maps here are the actual coefficient maps and
subgroup restrictions.  No other section normalizer
has been introduced into this relative pair.

\emph{The exact duality and half-rank comparison.}
The pair $(F,L)$ in (D.75.8) is the full completion
of a compact oriented genus-one surface with its
one boundary.  Apply the
surface-pair duality of
\cite[Propositions~5.7 and~6.16 of the accepted version]
{C-WilkesRelativeAccepted2018} to this specified pair.
Equivalently the dense discrete subgroup and its
cyclic boundary give the absolute and relative
cohomology comparisons, with their actual cup products,
as in (D.74.15).  For
$A=H^0(F,M)$ and $B=H^1(F,M)$, they give
\[
 H^0(F,\mathscr H_{M^\vee})=B^\vee,\qquad
 H^1(F,\mathscr H_{M^\vee})=A^\vee .
                                                               \tag{D.75.12}
\]
The identifications are equivariant for the actual
$E/F$ action.  The coefficient evaluation map is
\[
 M\otimes\mathscr H_{M^\vee}\longrightarrow\mathscr H_k,
 \qquad m\otimes\varphi\longmapsto
                 [d\longmapsto\varphi(d)(m)] .
                                                               \tag{D.75.13}
\]
Its target relative orientation group
$H^1(F,\mathscr H_k)=H^2(F,L;k)=k$ has trivial action.
In boundary coordinates this also follows from
$v c v^{-1}=c^{-1}$, which fixes the mod-two
boundary character.  All higher fibre rows vanish
because $\operatorname{cd}_2(F)=1$.

The two ordinary Hochschild--Serre sequences for
$M$ and $\mathscr H_{M^\vee}$ therefore have exactly
the rows in (D.74.21), with these $A$ and $B$.
Their products are induced by (D.75.13).
The positive-column cyclic complex has differential
$1+\sigma$, with the perfect kernel/image pairing
and period-one $u$ map of (D.74.18)--(D.74.19).
The product of the two fibre-row-one classes lies
in $H^2(F,\mathscr H_k)=0$.  Thus the same explicit
$d_2$ derivation identity gives
\[
                    \langle db,a^\vee\rangle
                          =\langle b,d'a^\vee\rangle .
                                                               \tag{D.75.14}
\]
These transgressions have equal rank, without
either one being assumed zero.

For a finite self-dual $M$, the high-degree
absolute and relative dimensions are consequently
equal.  Multiplication by the actual $u$ is an
isomorphism in sufficiently high degree on all
terms of (D.75.11), by the same two-row calculation.
Taking degree-zero localizations, write
\[
 \mathcal A=H^\bullet(E,M)[u^{-1}]_0,\quad
 \mathcal B=H^\bullet(D,M)[u^{-1}]_0,\quad
 \mathcal R=R^\bullet(E,M)[u^{-1}]_0 .
                                                               \tag{D.75.15}
\]
They are finite-dimensional, and the original
long exact sequence becomes
$\mathcal A\to\mathcal B\to\mathcal R\to\mathcal A$,
with the connecting map divided by $u$.
Its dimensions satisfy $\dim\mathcal A=\dim\mathcal R$.
The rank identities (D.74.27), applied to these
actual three maps, give
\[
                     \operatorname{rank}\bar r
                                  =\tfrac12\dim\mathcal B .
                                                               \tag{D.75.16}
\]
This use of D.74's coefficient argument has
replaced its three-cusp augmentation module
by (D.75.7); it has not replaced the torus
peripheral structure or its involution matrix.

\emph{Both incidence transports and the other section factor.}
Let $C=C_F(v)$.  Before proving it trivial,
the exact normalizer is
$N_E(S)=C\times S$, via multiplication and the
inverse $e\mapsto(ev^{-u(e)},v^{u(e)})$.
In the identity-modulo-two case let $T$ be the
other full section class of D.58 and write
$N_E(T)=C_F(T)\times T$ with the same labelled
coordinate maps.  In the other case there is
no such factor.

The finite-coefficient Brown theorem applies to
$E$ and to $D$: their open subgroups $F$ and
$L$ have cohomological dimension one, and every
finite subgroup injects into $C_2$.  Their high
cohomology is finite by the two-row sequences
with finite coefficients.  Thus the actual
dense restriction embeddings of
\cite[Theorem~2 and Section~1.1]
{A-KochloukovaZalesskii2007} are isomorphisms.
Put
\[
 U=H^\bullet(N_E(S),M)[u^{-1}]_0,\qquad
 V=H^\bullet(S,M)[u^{-1}]_0 .
                                                               \tag{D.75.17}
\]
Brown identifies $\mathcal A$ with $U$ together
with the other normalizer factor when present.
It identifies $\mathcal B$ with the two section
factors for $S,S'$.

Transport the first local factor by the identity,
and the second by the actual $g$ of (D.75.3).
The second cochain transport is
\[
 b(t_1,\ldots,t_n)\longmapsto
             g^{-1}b(gt_1g^{-1},\ldots,gt_ng^{-1}).
                                                               \tag{D.75.18}
\]
As proved by the explicit bar prism (D.74.33),
inner conjugation with exactly this coefficient
transport acts identically on global cohomology.
Consequently, if $r_S:U\to V$ is literal
restriction, the boundary map is exactly
\[
        \bar r:U\oplus U_{\rm other}\longrightarrow V\oplus V,
                 \qquad (a,a_{\rm other})
                             \longmapsto(r_Sa,r_Sa).
                                                               \tag{D.75.19}
\]
Here $U_{\rm other}=0$ in the one-class case.
The additional global section factor has no
boundary occurrence and its map in this
display is zero.  Both local sections occur
in the cusp class, by D.58's actual $F$-fusion.
Thus (D.75.16) says
$\operatorname{rank}r_S=\dim V$, proving
surjectivity of $r_S$ for every finite
self-dual coefficient module.

Suppose $1\ne k_0\in C$.  Choose an open normal
$N\triangleleft E$, with $N\subset F$ and
$k_0\notin N$, and use the exact finite module
\[
              J=E/N,\qquad \bar S=SN/N,\qquad
                             M=k[J/\bar S].
                                                               \tag{D.75.20}
\]
Its basis pairing is invariant and perfect.
In positive $S$-cohomology the singleton
$S$-fixed cosets are a basis; two-point orbits
have equal invariant and norm subspaces and
contribute zero.  The cosets $\bar S$ and
$\bar k_0\bar S$ are distinct and both fixed.
They are distinct because $\bar k_0\in F/N$
is nonidentity and $(F/N)\cap\bar S=1$.
The class of $\delta_{\bar S}$ is therefore
moved by $k_0$.

Every class in the image of $r_S$ is
$C$-invariant: the cochain map
$b\mapsto k_0 b(k_0^{-1}(\,\cdot\,)k_0)$ is
inner on $N_E(S)$, whereas on $S$ it is exactly
the coefficient action of $k_0$.
This contradicts the proved surjectivity.
Therefore $C=1$, proving (D.75.2).
All cusp reflections are $F$-conjugate to $v$:
inside $D$ each is conjugate to $v$ or $cv$,
and these two fuse by $g$.  Their fixed
subgroups in $F$ are therefore all trivial.
Dividing any two solutions of (D.75.3) gives
an element of $C$, so that solution is unique.

\emph{Every unit and both literal cusp reflections.}
The two half-exponents in (D.75.4) belong
to $\widehat{\mathbb Z}$, because a full unit
is odd at two and division by two is defined
at every other prime.  The identities
\[
 c^a v c^{-a}=c^{2a}v,\qquad
                 gvg^{-1}=cv,\qquad g^{-1}(cv)g=v
                                                               \tag{D.75.21}
\]
show directly that
$a_{\mu,i}v_i a_{\mu,i}^{-1}=c^\mu v_i$.
Dividing two such transporters gives an
element of $C_F(v_i)=1$.  This proves
the full singleton assertion (D.75.4).
Since $c$ is a commutator, D.59's primitive
homology statement gives
\[
 [a_{\mu,i}]_{\rm ab}=(-1)^i[g]_{\rm ab},\qquad
          \alpha_i(a_{\mu,i})=c^{-\mu}a_{\mu,i},
      \quad \alpha_i=\operatorname{Inn}(v_i)|_F .
                                                               \tag{D.75.22}
\]
In particular D.62's full primitive-homology
centralizer theorem applies to every one of
these elements, without a free-basis
primitivity assumption.

For completeness retain the maps relating
(D.75.5) to the peripheral conjugator.
Fix $(\mu,i)$ and write $a=a_{\mu,i}$.
Every outside element is uniquely $hv_i$,
$h\in F$, and
\[
             (hv_i)a(hv_i)^{-1}
                         =h(c^{-\mu}a)h^{-1}.
                                                               \tag{D.75.23}
\]
Thus the required $F$-conjugacy is exactly
the existence of an outside centralizer
element.  Its kernel under $u$ is the full
group $\overline{\langle a\rangle}\cong
\widehat{\mathbb Z}$ by D.62.
Choose an outside centralizer element $t$;
then $t^2=a^\kappa$, and the complete
centralizer has the actual coordinates
$a^r t^\epsilon$, as in (D.62k)--(D.62l).

In the two-class case the D.60 character
$\chi:E\to\widehat{\mathbb Z}$ has
$\chi(a)=(-1)^i$.  Hence
$2\chi(t)=(-1)^i\kappa$ and $\kappa_2$ is even.
There is the unique outside commuting involution
\[
                         w=a^{-\kappa/2}t .
                                                               \tag{D.75.24}
\]
It belongs to the other class, because a
cusp-class involution has trivial fixed
subgroup in $F$ and $a\ne1$.

In the one-class case every outside involution
has trivial fixed subgroup in $F$, so no such
involution centralizes $a\ne1$.  Thus
$\kappa_2$ is odd.  The full root is obtained
by the exact primewise correction
\[
 \begin{gathered}
 b_2=0,\qquad b_p=(1-\kappa_p)/2\quad(p\ne2),\\
 \kappa'=\kappa+2b\in\widehat{\mathbb Z}^{\times},
 \qquad z=(a^b t)^{(\kappa')^{-1}},
 \qquad z^2=a .
 \end{gathered}                                             \tag{D.75.25}
\]
The centralizer is the full procyclic group
$\overline{\langle z\rangle}$; multiplication
by two in this group is injective, so this
root is unique.  These are the full-coordinate
maps of D.62, rather than a root only in a
maximal pro-two quotient.

The resulting extension maps from the
specified D.61 standard models are
\[
 \begin{array}{c|ccc}
       &X&Y&v_0\\ \hline
 \text{two classes}&a&v_iw&v_i\\
 \text{one class}&zv_i&v_i z&v_i .
 \end{array}
 \qquad [X,Y]\longmapsto c^\mu .
                                                               \tag{D.75.26}
\]
The commutator identity in the first row
uses $[a,w]=1$; in the second it uses
$z^2=a$ and (D.75.4).  Conjugation by
$v_i$ has exactly the standard D.61 action
on the two displayed basis images.
D.61's actual boundary test makes the kernel
map a full isomorphism.  The inverse
extension map is its inverse on $F$ and
sends the specified $v_i$ back to $v_0$.
Conversely every D.61 normalized
peripheral conjugator has a unit $\mu$
and one of these labels $i$.  Its image
of $X$ in the two-class model, or of $XY$
in the one-class model, is the unique
element in (D.75.4).  The image of its
other reflection or of $Xv_0$, respectively,
gives the outside centralizer element in
(D.75.23).  This proves (D.75.5) in both
directions with all cusp and quotient labels.

\emph{The one-class consequence and the other fixed group.}
If $\alpha_*\bmod2\ne I$, D.58 gives only
the cusp section class.  Every involution
is therefore self-centralizing in $E$.
The actual open subgroup $F$ is projective
of index two, and the involutions form
the closed subset
$\{t:t^2=1\}\cap(E\setminus F)$.
Haran's Proposition~2.2
\cite{C-HaranArtinSchreier1993} applies.
The resulting real-projectivity has
exactly the finite embedding-problem maps
spelled out in (D.74.38).

If $\alpha_*\bmod2=I$, D.58 gives the other
full section class $T=\langle w_1\rangle$.
D.59 retains the exact assertions
\[
 \begin{gathered}
 N_E(T)=C_F(w_1)\times T,\qquad
             \dim_k H^1(C_F(w_1),k)=1,\\
 C_F(w_1)\longrightarrow
       \ker(\alpha_*-I:H_1(F,\widehat{\mathbb Z})
                            \to H_1(F,\widehat{\mathbb Z}))
                  \quad\text{is onto}.
 \end{gathered}                                             \tag{D.75.27}
\]
This factor is precisely the unused factor
in (D.75.19).  The proof does not identify
it with a procyclic group.  In either case,
the remaining recognition question is the
exact conjugacy in (D.75.5); its existence
has not been inferred from the uniqueness
of the transporter.
\end{proof}
